\subsubsection{Mellin 尺度递归系数、双基数去混叠与递归证书复杂度下界}\label{subsubsec:xi-mellin-scale-recursion-bohr-dealiasing}
本段把“零点的尺度递归表征”写成可直接审计的 Mellin 系数链条，并把“平滑/截断诱导的误差控制”与“递归定位位深下界”放在同一闭环：前者给出 RH 的尺度有界性表述与双基数去混叠编码，后者给出紧化坐标下不可规避的 dyadic 分辨率成本。

\paragraph{尺度递归系数与 Mellin 归一化}
记
\[
\psi(x):=\sum_{n\le x}\Lambda(n),\qquad
\psi_0(x):=\frac{\psi(x+0)+\psi(x-0)}{2},
\]
并定义归一化偏差
\[
E(x):=\frac{\psi_0(x)-x}{x^{1/2}}\qquad (x>1).
\]
固定伸缩基数 \(\lambda>1\)。对任意 \(\varphi\in C_c^\infty((1,\lambda))\) 且满足零 Mellin 均值
\[
\mathcal M\varphi(0)=\int_0^\infty \varphi(u)\,\frac{\dd u}{u}=0,
\]
定义尺度递归系数
\[
C_{\varphi,\lambda}(m):=
\int_1^\infty E(x)\,\varphi\!\left(\frac{x}{\lambda^m}\right)\frac{\dd x}{x},
\qquad m\in\NN.
\]
其中 Mellin 变换记为
\[
\mathcal M\varphi(s):=\int_0^\infty \varphi(u)\,u^s\frac{\dd u}{u}.
\]

称有限族 \(\Phi=\{\varphi_r\}_{r=1}^R\subset C_c^\infty((1,\lambda))\) 为非退化，若
\begin{enumerate}
  \item 对所有 \(r\) 有 \(\mathcal M\varphi_r(0)=0\)；
  \item 对任意非平凡零点 \(\rho\) 满足 \(\Re\rho\neq \frac12\) 时，存在 \(r\) 使
  \(\mathcal M\varphi_r(\rho-\frac12)\neq 0\)。
\end{enumerate}

\begin{theorem}[Bohr--Mellin 多尺度判据：RH 的尺度有界/几乎周期等价]\label{thm:xi-mellin-scale-bohr-rh-equivalence}
给定 \(\lambda>1\) 与一个非退化族 \(\Phi=\{\varphi_r\}_{r=1}^R\)，下列陈述等价：
\begin{enumerate}
  \item RH 成立；
  \item 对每个 \(r\)，序列 \(\{C_{\varphi_r,\lambda}(m)\}_{m\ge 1}\) 一致有界；
  \item 对每个 \(r\)，存在 Bohr 几乎周期函数 \(F_r:\RR\to\CC\) 使
  \[
  C_{\varphi_r,\lambda}(m)=F_r(m)+O(\lambda^{-m/2}),
  \]
  且 \(F_r\) 的 Bohr 频谱由集合
  \(\{\gamma\log\lambda:\ \rho=\frac12+\mathrm{i}\gamma\}\) 生成。
\end{enumerate}
\end{theorem}
\begin{proof}
由 \(\psi_0\) 的显式公式（见 \cite{Titchmarsh1986RiemannZeta}）可写
\[
E(x)=
-\sum_{\rho}\frac{x^{\rho-\frac12}}{\rho}
+O(x^{-1/2})\qquad (x\ge 2),
\]
其中求和遍历非平凡零点并按重数计。代入 \(C_{\varphi,\lambda}(m)\) 后，令 \(x=\lambda^m u\)，得到
\begin{equation}\label{eq:xi-mellin-scale-coeff-explicit}
C_{\varphi,\lambda}(m)
=-\sum_{\rho}
\frac{\lambda^{m(\rho-\frac12)}\mathcal M\varphi(\rho-\frac12)}{\rho}
+O(\lambda^{-m/2}).
\end{equation}
交换求和与积分的合法性来自两点：其一，\(\varphi\in C_c^\infty\) 蕴含
\(\mathcal M\varphi(\sigma+\mathrm{i}t)\) 沿竖线对任意阶快速衰减；其二，零点计数满足
\(N(T)=O(T\log T)\)，故与快速衰减权相乘后绝对可和。

若 RH 成立，\(\rho=\frac12+\mathrm{i}\gamma\)，于是
\[
C_{\varphi,\lambda}(m)
=-\sum_{\gamma}
\frac{\mathcal M\varphi(\mathrm{i}\gamma)}{\frac12+\mathrm{i}\gamma}
e^{\mathrm{i}m\gamma\log\lambda}
+O(\lambda^{-m/2}).
\]
主和是绝对收敛的指数级数，因此给出 Bohr 几乎周期函数 \(F_\varphi(m)\)；
进而得到 (1)\(\Rightarrow\)(3)\(\Rightarrow\)(2)。

反过来，若 RH 不成立，则存在 \(\rho_0=\beta_0+\mathrm{i}\gamma_0\) 且 \(\beta_0>\frac12\)。
由非退化性可取某个 \(r\) 使
\(\mathcal M\varphi_r(\rho_0-\frac12)\neq 0\)。
在 \eqref{eq:xi-mellin-scale-coeff-explicit} 中，相应模态携带因子
\(\lambda^{m(\beta_0-\frac12)}\)，其幅度随 \(m\) 指数放大；该增长不可能与余下模态及 \(O(\lambda^{-m/2})\) 项共同维持全局一致有界，故 (2) 不成立。
因此 (2)\(\Rightarrow\)(1)。
\end{proof}

\begin{theorem}[双基数去混叠的相位注入编码]\label{thm:xi-dual-base-dealiasing-injective-encoding}
\leanverified{paper\_xi\_dual\_base\_dealiasing\_injective\_encoding}
设 RH 成立，取 \(\lambda_1,\lambda_2>1\) 且
\begin{equation}\label{eq:xi-dual-base-irrational-log-ratio}
\frac{\log\lambda_1}{\log\lambda_2}\notin\QQ.
\end{equation}
则映射
\[
\gamma\longmapsto
\left(e^{\mathrm{i}\gamma\log\lambda_1},\ e^{\mathrm{i}\gamma\log\lambda_2}\right)\in\TT^2
\]
在零点纵坐标集合上为单射。换言之，若
\[
\left(e^{\mathrm{i}\gamma\log\lambda_1},e^{\mathrm{i}\gamma\log\lambda_2}\right)
=
\left(e^{\mathrm{i}\gamma'\log\lambda_1},e^{\mathrm{i}\gamma'\log\lambda_2}\right),
\]
则必有 \(\gamma=\gamma'\)。
\end{theorem}
\begin{proof}
若上述相位对相等，则
\[
(\gamma-\gamma')\log\lambda_j\in 2\pi\ZZ,\qquad j=1,2.
\]
于是存在 \(k_1,k_2\in\ZZ\) 使
\[
\gamma-\gamma'
=\frac{2\pi k_1}{\log\lambda_1}
=\frac{2\pi k_2}{\log\lambda_2}.
\]
若 \(\gamma\neq\gamma'\)，则 \((k_1,k_2)\neq(0,0)\)，从而
\(
\log\lambda_1/\log\lambda_2=k_1/k_2\in\QQ
\)，与 \eqref{eq:xi-dual-base-irrational-log-ratio} 矛盾。故 \(\gamma=\gamma'\)。
\end{proof}

\begin{corollary}[二进制--黄金比去混叠对]\label{cor:xi-dual-base-dealiasing-2-phi}
\leanverified{paper\_xi\_dual\_base\_dealiasing\_2\_phi}
取 \(\lambda_1=2\)、\(\lambda_2=\varphi=\frac{1+\sqrt5}{2}\)，则定理
\ref{thm:xi-dual-base-dealiasing-injective-encoding} 的单射结论成立。
\end{corollary}
\begin{proof}
若 \(\log\varphi/\log2\in\QQ\)，则存在 \(p,q\in\NN\) 使 \(\varphi^q=2^p\)。
在二次域 \(\QQ(\sqrt5)\) 取范数：
\[
(-1)^q=\mathrm N(\varphi^q)=\mathrm N(2^p)=4^p,
\]
矛盾。故 \(\log\varphi/\log2\notin\QQ\)。
\end{proof}

\begin{theorem}[零点谱截断的统一超多项式误差界]\label{thm:xi-scale-spectral-truncation-uniform-superpoly}
\leanverified{paper\_xi\_scale\_spectral\_truncation\_uniform\_superpoly}
在 RH 下，固定 \(\lambda>1\) 与 \(\varphi\in C_c^\infty((1,\lambda))\)、\(\mathcal M\varphi(0)=0\)。
由 \eqref{eq:xi-mellin-scale-coeff-explicit} 写
\[
C_{\varphi,\lambda}(m)
=-\sum_{\gamma}
\frac{\lambda^{\mathrm{i}m\gamma}\mathcal M\varphi(\mathrm{i}\gamma)}{\frac12+\mathrm{i}\gamma}
+R_{\varphi,\lambda}(m),
\qquad
\left|R_{\varphi,\lambda}(m)\right|\le C_{\varphi,\lambda}\lambda^{-m/2}.
\]
定义截断量
\[
C_{\varphi,\lambda}^{(\Gamma)}(m):=
-\sum_{|\gamma|\le\Gamma}
\frac{\lambda^{\mathrm{i}m\gamma}\mathcal M\varphi(\mathrm{i}\gamma)}{\frac12+\mathrm{i}\gamma}
+R_{\varphi,\lambda}(m).
\]
则对任意 \(A>0\)，存在常数 \(K_A(\varphi,\lambda)\) 使
\[
\sup_{m\ge 1}
\left|C_{\varphi,\lambda}(m)-C_{\varphi,\lambda}^{(\Gamma)}(m)\right|
\le
K_A(\varphi,\lambda)\,\Gamma^{-A},
\qquad \Gamma\ge 2.
\]
\end{theorem}
\begin{proof}
两者差仅为零点尾和：
\[
\left|C_{\varphi,\lambda}(m)-C_{\varphi,\lambda}^{(\Gamma)}(m)\right|
\le
\sum_{|\gamma|>\Gamma}
\frac{\left|\mathcal M\varphi(\mathrm{i}\gamma)\right|}{\left|\frac12+\mathrm{i}\gamma\right|}.
\]
对任意 \(B>0\)，有
\(\left|\mathcal M\varphi(\mathrm{i}\gamma)\right|\ll_B(1+|\gamma|)^{-B}\)；
再结合 \(N(T)=O(T\log T)\) 给出的单位长度零点数估计，得
\[
\sum_{|\gamma|>\Gamma}
\frac{\left|\mathcal M\varphi(\mathrm{i}\gamma)\right|}{\left|\frac12+\mathrm{i}\gamma\right|}
\ll_B
\sum_{n\ge \Gamma}\frac{\log n}{n^{B+1}}
\ll_A \Gamma^{-A},
\]
取 \(B\) 足够大即得结论。
\end{proof}

\begin{theorem}[紧化坐标下的 dyadic 最小位深下界]\label{thm:xi-compactified-dyadic-depth-lower-bound}
\leanverified{paper\_xi\_compactified\_dyadic\_depth\_lower\_bound}
令 \(0<\gamma_1\le\gamma_2\le\cdots\) 为非平凡零点纵坐标（按重数计），并定义
\[
u(t):=\frac{1}{\pi}\arctan t\in\left(0,\frac12\right),\qquad
u_n:=u(\gamma_n).
\]
则当 \(T\to\infty\) 时，在区间 \([T,2T]\) 内存在相邻零点
\(\gamma_n<\gamma_{n+1}\) 满足
\begin{equation}\label{eq:xi-u-gap-upper-bound}
u_{n+1}-u_n
\le
\frac{2}{(1+T^2)\log\!\bigl(\frac{T}{2\pi}\bigr)}\,(1+o(1)).
\end{equation}
因此，若某 dyadic 递归隔离方案要求在 \(u\)-坐标上使用长度至多 \(2^{-b}\) 的区间并满足“每区间至多含一个零点”，则必有
\begin{equation}\label{eq:xi-dyadic-bitdepth-lower-bound}
b\ge
\log_2\!\Bigl((1+T^2)\log\!\bigl(\tfrac{T}{2\pi}\bigr)\Bigr)-1+o(1)
=2\log_2 T+\log_2\log T+O(1).
\end{equation}
\end{theorem}
\begin{proof}
由 Riemann--von Mangoldt 公式（见 \cite{Titchmarsh1986RiemannZeta}）
\[
N(T)=\frac{T}{2\pi}\log\frac{T}{2\pi}-\frac{T}{2\pi}+O(\log T),
\]
可得
\[
N(2T)-N(T)=\frac{T}{2\pi}\log\frac{T}{2\pi}+O(T).
\]
故 \([T,2T]\) 内平均零点间距为
\[
\frac{T}{N(2T)-N(T)}
=\frac{2\pi}{\log\!\bigl(\frac{T}{2\pi}\bigr)}(1+o(1)),
\]
从而存在相邻零点满足
\[
\gamma_{n+1}-\gamma_n
\le
\frac{2\pi}{\log\!\bigl(\frac{T}{2\pi}\bigr)}(1+o(1)).
\]
又 \(u'(t)=1/(\pi(1+t^2))\)，在 \([T,2T]\) 上有
\[
u'(\xi)\le \frac{1}{\pi(1+T^2)}.
\]
由中值定理即得 \eqref{eq:xi-u-gap-upper-bound}。
若 \(2^{-b}\) 大于某个相邻差值，则无法保证“每区间至多一个零点”，故必须有
\(
2^{-b}\le \min(u_{n+1}-u_n)
\)
于目标层级，代入 \eqref{eq:xi-u-gap-upper-bound} 得 \eqref{eq:xi-dyadic-bitdepth-lower-bound}。
\end{proof}

\begin{corollary}[递归零点证书的总比特复杂度下界]\label{cor:xi-recursive-certificate-bit-complexity-lower-bound}
记 \(B(T)\) 为在 \(u\)-坐标中对所有 \(0<\gamma_n\le T\) 构造两两不重叠 dyadic 证书区间所需的最小总位深预算。则
\[
B(T)\ge
\bigl(N(T)-N(T/2)\bigr)\,
\bigl(2\log_2T+\log_2\log T+O(1)\bigr),
\]
从而
\[
B(T)=\Omega\!\bigl(T(\log T)^2\bigr).
\]
\end{corollary}
\begin{proof}
对高度层 \([T/2,T]\) 内的每个零点，定理 \ref{thm:xi-compactified-dyadic-depth-lower-bound}
给出同阶位深下界 \(2\log_2T+\log_2\log T+O(1)\)。
该层零点个数为 \(N(T)-N(T/2)\asymp T\log T\)，将逐点下界求和即得结论。
\end{proof}

\endinput
